\subsection{The exact remaining pentagon value}

\begin{statement}{D23.P25}{Proved}{The required value at the fourth--fifth pair is a commutator}
\end{statement}
Use the positive braid marking
\[
 X_1=A^2,\quad X_2=B^2,\quad X_3=C^2,\quad
 X_4=(ABA)^2,\quad X_5=(CBC)^2,
\]
and put $D=ABCABA$. Let $f$ satisfy the two-cycle, and
write $F=f(X_1,X_2)$ and $R=f(X_3,X_4)$. In the
five-factor pentagon order of supplied FP.37.3, the
pentagon equation is equivalent to
\[
 \boxed{\quad f(X_4,X_5)=[D,FR]
            =DFRD^{-1}R^{-1}F^{-1}.\quad}                 \tag{D23.P25.1}
\]
Here $[a,b]=aba^{-1}b^{-1}$. The formula keeps the actual
inverse of $D$ and remains valid in the raw braid action,
where $D^2$ is the nontrivial full-twist deck generator.

\begin{proof}
The reversing half twist satisfies
\[
 DX_1D^{-1}=X_3,\quad DX_2D^{-1}=X_2,\quad
 DX_4D^{-1}=X_5.
\]
It also sends $X_3$ back to $X_1$, since $D^2$ is
central. Equivariance of evaluation and the two-cycle give
\[
 f(X_5,X_1)=DR^{-1}D^{-1},\qquad
 f(X_2,X_3)=DF^{-1}D^{-1}.
\]
Thus, putting $H=f(X_4,X_5)$, the literal pentagon is
\[
 \mathcal P(f)
   =R(DR^{-1}D^{-1})(DF^{-1}D^{-1})HF
   =RDR^{-1}F^{-1}D^{-1}HF.
\]
Solving $\mathcal P(f)=1$ for $H$ gives
$H=DFRD^{-1}R^{-1}F^{-1}$, exactly (D23.P25.1).
No involutivity assumption on $D$ was used.
\end{proof}

\paragraph{Exact replay and remaining test.}
\path{hurwitz_required_f45.py} verifies all displayed
conjugacies in both the 1344-point and 8064-point actions,
and verifies $D^2$ equals the literal full twist. For
all 36 recorded D23.H7 witnesses it computes the
required value in (D23.P25.1) and checks the original
five-factor product entry by entry. The raw required
value projects to the base required value in every case.
The NPZ stores all these required values.

These identities do not prove that the required value
belongs to $\langle X_4,X_5\rangle'$, or that a common
word realizes $F,R$ and that value while satisfying
the gentle identities. They give the exact remaining
value to impose in those tests.
